\section{Residual geometry}
\begin{definition}[Normalized residual geometry]\label{def:normalized-residual-geometry}
The normalized residual space carries:
\begin{enumerate}[leftmargin=2em]
  \item admissible descent sectors and strict descent sectors,
  \item trapping close basins,
  \item honest promote regions,
  \item a residual frontier separating the geometrically unresolved zone.
\end{enumerate}
\end{definition}

\begin{definition}[Residual frontier]\label{def:residual-frontier}
The residual frontier is the part of normalized residual space lying in neither a trapping close basin nor an honest promote region.
\end{definition}

\begin{theorem}[Geometric control theorem]\label{thm:geometric-control}
In a fixed canonical normalization regime, the normalized residual space admits a route-faithful geometric control structure consisting of admissible descent sectors, close basins, honest promote regions, a residual frontier, and descent selectors on finite admissible successor families.
\end{theorem}

\begin{proof}
These structures are defined entirely within the normalized residual space and remain inside the same route-faithful invariant class. Each construction therefore preserves the inherited horizon/boundary distinction, earliest-boundary stop, and certificate-neutrality.
\end{proof}

